\subsection{Schopnost pojmout celé studentské řešení}
\label{subsec:criteria-context-window}

V \secref[sekci]{subsec:llm-token-limit} byla popsána omezená kapacita kontextového okna LLM, která určuje, kolik tokenů může model zpracovat v jednom požadavku.
Pro analýzu studentských řešení je klíčové, aby model dokázal pojmout celé odevzdání včetně doprovodného kontextu v jednom požadavku, protože pouze tak může poskytnout kontextově přesnou zpětnou vazbu.
Vstup modelu se typicky skládá z následujících částí:

\begin{itemize}
    \item \textbf{Systémový prompt} -- instrukce definující roli modelu, požadavky na formát výstupu a další parametry generování.
    Typicky 200--500 tokenů.

    \item \textbf{Zadání úlohy} -- popis úlohy včetně požadavků na vstupy, výstupy a případné specifické podmínky.
    Může být poměrně rozsáhlé, zejména pokud obsahuje příklady vstupů a výstupů -- přibližně 500--1~000 tokenů.

    \item \textbf{Zdrojový kód} -- samotné řešení studenta.
    Pro středně velkou úlohu se jedná o 1~000--5~000 tokenů.
\end{itemize}

Pro standardní studentské odevzdání v systému Kelvin je tedy dostačující kontextové okno přibližně \textbf{8~000 tokenů}.
S ohledem na rozsáhlejší úlohy (větší projekty, vícesouborová řešení) a princip dostatečné rezervy je jako \textbf{minimální požadavek} stanoven rozsah \textbf{16~000 tokenů}.
Moderní modely tuto hranici zpravidla bez problémů splňují, jelikož většina aktuálně dostupných modelů nabízí kontextové okno v rozsahu 32~000 až 128~000 tokenů.

\begin{table}[H]
    \centering
    \resizebox{\textwidth}{!}{%
        \begin{tabular}{lp{9cm}l}
            \toprule
            \textbf{Kritérium}              & \textbf{Popis}                                                                 & \textbf{Minimální požadavek} \\
            \midrule
            Kvalita výstupů                 & Věcně správné, srozumitelné, relevantní komentáře                              & Vyšší než průměr \\
            Dostupnost                      & Přístupný přes API nebo ke stažení; vhodná licence pro akademické použití      & Open-source nebo API s DPA \\
            Náročnost provozu               & Přijatelné náklady na token nebo na hardware                                   & Viz \tabref{tab:deployment-comparison} \\
            Kontextové okno                 & Schopnost zpracovat celé odevzdání včetně promptu                              & $\geq$ 16~000 tokenů \\
            \bottomrule
        \end{tabular}
    }
    \caption{Přehled kritérií výběru modelu a jejich minimálních požadavků}
    \label{tab:selection-criteria}
\end{table}

\endinput